\section{Introduction}
Human pose estimation refers to the task of locating the principal body
joints—such as elbows, knees, and shoulders—within an image. The exact list of
joints depends on the data set: \textsc{COCO-WholeBody} annotates seventeen
body keypoints in addition to detailed hand, foot, and face landmarks
\cite{jin2020whole}, whereas the Leeds Sports Pose data set provides fourteen
body keypoints \cite{johnson2010clustered}. Over the past decade the field has
converged on two broad pipelines. In a \emph{top-down} pipeline, each person is
first enclosed in a bounding box and the keypoints are predicted inside that
region; in a \emph{bottom-up} pipeline, every visible joint is detected in a
single pass and a grouping stage later assembles them into individual
skeletons. Before the deep-learning era, solutions combined handcrafted image
descriptors with relatively shallow learners. Examples include relevance-vector
machines \cite{tipping2001sparse}, mixtures of experts
\cite{sigal2006predicting}, locality-constrained coding
\cite{wang2010locality}, relevant-component analysis \cite{kanaujia2007semi},
bag-of-words representations \cite{ning2008discriminative}, random forests
\cite{chang2013fast}, Hough forests \cite{girshick2011efficient}, pixel- and
superpixel-level models \cite{wang2011multi}, and motion cues
\cite{sapp2013modec}. While these methods laid the essential groundwork,
limited representational capacity hampered their performance and they naturally
struggled to compete with more modern approaches.

Deep convolutional networks (CNNs) \cite{lecun1998gradient} transformed the
task by learning features directly from data. \textsc{DeepPose} (revolutionary
at the time) \cite{toshev2014deeppose}, based on the eight-layer AlexNet
backbone \cite{krizhevsky2012imagenet}, was the first model to show that an
end-to-end CNN could localise human joints reliably. Subsequent architectures
escalated both depth and representational capacity: VGGnet with 16-19 layers
\cite{simonyan2014very}; ResNet with 50, 101, or 152 layers linked by residual
shortcuts \cite{he2016deep}; and, later, High-Resolution Networks (HRNet) that
maintain parallel feature streams at multiple resolutions throughout the
network \cite{sun2019deep}. On top of these backbones, researchers introduced
multi-stage refinement modules that progressively sharpen coarse heat-maps into
precise keypoint locations—examples include Convolutional Pose Machines
\cite{wei2016convolutional}, stacked hourglass networks
\cite{newell2016stacked}, and the part-affinity field branch used by
\textsc{OpenPose} \cite{cao2017realtime}. Additional variants (all using some
convolution backbone), such as Cascade Feature Aggregation
\cite{su2019cascade}, SimpleBaseline \cite{xiao2018simple}, and the Cascaded
Pyramid Network \cite{chen2018cascaded}, improved accuracy, but each gain was
accompanied by additional layers, parameters, and computational cost (this is a
relatively broad statement but it does describe the trend).

More recently, vision transformers have emerged as a versatile alternative to
convolutional backbones, replacing local convolution with multi-head
self-attention so that every image patch -- or any intermediate "token" standing
in for a joint -- can exchange information with all others in a single layer
\cite{vaswani2017attention}. In practice the image (or a coarse set of
heat-maps) is converted into a token sequence, processed by stacked transformer
blocks that model long-range dependencies, and then decoded back into keypoint
coordinates. Variations on this theme include \textsc{TokenPose}, which
attaches a learnable token to each joint and passes the sequence through a
transformer encoder \cite{li2021tokenpose}; \textsc{HRFormer}, which retains
HRNet's multi-resolution branches but fuses them with windowed attention
\cite{yuan2021hrformer}; and \textsc{ViTPose}, which drops the convolutional
stem entirely and shows that a plain Vision Transformer can match or surpass
CNN baselines when enough data are available \cite{xu2022vitpose}. Transformer
backbones are offered in graded sizes -- small models that run in real time on
consumer GPUs, medium models that balance speed and accuracy, and large models
whose parameter counts approach a billion or more -- so speed/accuracy trade-offs
can be made by selecting a lighter or heavier version of the \emph{same
architecture} rather than rebuilding the entire network.

Depth cues provide a complementary signal that helps disambiguate overlapping
limbs. Monocular predictors such as \textsc{MiDaS}, trained on a broad
collection of real-world data sets \cite{ranftl2021robust,ranftl2022dense}, and
\textsc{ZoeDepth}, which mixes real and synthetic images to widen scene
coverage \cite{khattak2023zoedepth}, now produce accurate, scale-consistent
depth maps with a single forward pass. Attaching the estimated depth value to
each keypoint effectively lifts the 2-D skeleton into 2.5-D space, offering
geometric context that pure image coordinates lack.

Multi-person scenes remain challenging because joints belonging to different
individuals often overlap or lie in close proximity, making \emph{pose
tracking} -- the task of grouping detected joints into the correct person-specific
skeletons -- non-trivial. One approach is to represent the estimated skeleton as a
graph in which each node corresponds to a joint and each edge encodes spatial
or geometric relationships between joints; a Graph Attention Network (GAT) then
learns attention weights that highlight the most informative neighbours
\cite{velivckovic2017graph}. This paper proposes a \emph{depth-aware} GAT that
combines 2-D image coordinates with per-joint depth values, remains
detector-agnostic -- able to consume keypoints from either a lightweight ResNet
heat-map head or a transformer front-end (as examples) -- and is designed to
improve the accuracy and real-time performance of pose tracking in crowded
scenes.